% A hand-written paper, not made by PaperTeX
\documentclass[12pt,letterpaper]{amsart}
\usepackage[utf8]{inputenc}
\usepackage{amsmath,amssymb,amsthm}
\usepackage[margin=1in]{geometry}
\usepackage{graphicx}
\usepackage{siunitx}
\usepackage{listings}
\usepackage[colorlinks=true]{hyperref}

\newtheorem{thm}{Theorem}[section]
\newtheorem{lem}[thm]{Lemma}
\newtheorem{conj}{Conjecture}
\theoremstyle{definition}
\newtheorem{defn}{Definition}

\newcommand{\R}{\mathbb{R}}
\DeclareMathOperator{\tr}{tr}

\title{On the Stability of Small Things\thanks{Supported by nobody.}}
\author{Ana Pérez \\ University of Somewhere \\ \texttt{ana@example.org} \and Bo Li}
\date{}

\begin{document}
\maketitle

\begin{abstract}
We show that small things are stable. % a comment that should vanish
This uses \emph{two} ideas.
\end{abstract}

\section{Introduction}\label{sec:intro}
Let $f\colon \R \to \R$ be smooth, and write $\|x\|$ for the norm. Following
Erd\H{o}s and Gödel---and also Poincar\'e~\cite[p.~3]{poincare,godel}---we
consider ``quoted'' text, {\bf bold}, {\em emphasis}, \textsc{Small Caps}, 50\% of
cases and \$5. See Section~\ref{sec:main} and \eqref{eq:main}.\footnote{A note with $x^2$ in it.}

\[ E = mc^2 \]
and also
$$ a^2 + b^2 = c^2 $$

\begin{align}
  f(x) &= \sum_{n=0}^{\infty} \frac{x^n}{n!} \label{eq:main} \\
       &= e^x \nonumber
\end{align}

\begin{equation*}
  \begin{split}
    g(x) &= \int_0^x t\,dt \\
         &= \frac{x^2}{2}
  \end{split}
\end{equation*}

\section{Main result}
\label{sec:main}

\begin{thm}[Stability]\label{thm:stab}
Every small thing is stable.
\end{thm}

\begin{proof}
Trivial, since
\[ 0 \le \tr A \le 1. \qedhere \]
\end{proof}

\begin{conj}
Large things are stable too.
\end{conj}

\begin{defn}
A thing is \emph{small} if it fits.
\end{defn}

\subsection{Lists}
\begin{enumerate}
  \item First
  \item Second, with
  \begin{enumerate}
    \item a nested one
    \item and another
  \end{enumerate}
  \item Third
\end{enumerate}

\begin{description}
  \item[Term] Its meaning.
  \item[Other] Another meaning.
\end{description}

\begin{table}[h]
\centering
\begin{tabular}{|l|c|r|}
\hline
\textbf{Name} & \textbf{Value} & \textbf{Unit} \\
\hline
Mass & 1.5 & kg \\
Speed & \multicolumn{2}{c|}{fast} \\
\hline
\end{tabular}
\caption{Some values.}
\label{tab:vals}
\end{table}

\begin{figure}[t]
  \centering
  \includegraphics[width=0.5\textwidth]{plot.pdf}
  \caption{A plot.}\label{fig:plot}
\end{figure}

\begin{lstlisting}[language=Python, caption={Hello}]
print("hi")  # 100% sure
\end{lstlisting}

\begin{verbatim}
raw % text { here
\end{verbatim}

\begin{itemize}
\item Odd \verb|code| inline
\end{itemize}

\vspace{1cm}
\SI{3}{\metre} is a length.

\input{appendix}

\begin{thebibliography}{9}
\bibitem{poincare} H.~Poincar\'e, \emph{Science and Hypothesis}, Walter Scott, 1905.
\bibitem{godel} K. Gödel. On formally undecidable propositions. 1931.
\end{thebibliography}

\end{document}
